\documentclass[11pt]{article}

% ============================================================
% Milestone 1 -- Project Proposal & Literature Review
% DIP Team Project -- Project 28: New No-Reference Image-Quality Metric
% Compile: pdflatex -> bibtex -> pdflatex -> pdflatex
% ============================================================
\usepackage[margin=1in]{geometry}
\usepackage[T1]{fontenc}
\usepackage[utf8]{inputenc}
\usepackage{amsmath,amssymb}
\usepackage{graphicx}
\usepackage{booktabs}
\usepackage[round]{natbib}
\usepackage{hyperref}
\usepackage{xurl}
\hypersetup{colorlinks=true, linkcolor=blue!50!black, citecolor=blue!50!black, urlcolor=blue!50!black}
\usepackage{enumitem}
\usepackage{xcolor}

\newcommand{\todo}[1]{\textcolor{red}{\textbf{[TODO: #1]}}}

\title{\LARGE Project Proposal \& Literature Review \\[4pt]
\large Building and Validating a Lightweight No-Reference Image-Quality Metric (Project 28)}
\author{
  \todo{Student A Full Name} \and
  \todo{Student B Full Name} \and
  \todo{Student C Full Name}
}
\date{Milestone 1 --- \todo{Date} \\{} [COURSE CODE] --- Digital Image Processing}

\begin{document}
\maketitle

\begin{abstract}
Automated image-quality assessment is essential wherever images are captured, compressed, or
enhanced, but the strongest perceptual metrics (PSNR, SSIM, LPIPS) all require a pristine
reference image that does not exist in most real deployments. We propose a lightweight
no-reference image-quality metric: a small regression head trained on features of a frozen
pretrained backbone, using the KADID-10k database (10,125 distorted images with human
opinion scores). Our feasibility study already reproduces the reference baselines on a
content-based split (seed 42): LPIPS achieves SROCC~0.808, SSIM~0.627, and PSNR~0.667
against human scores, giving the target our no-reference model must approach without ever
seeing the reference. The code repository and results are public at
\url{https://github.com/mo1998/dip-project-28-nr-iqa}.
\end{abstract}

\section{Problem Statement \& Motivation}
\label{sec:motivation}
Given a single image that has been distorted (compressed, blurred, noised, re-touched) and
\emph{no} pristine original to compare against, we must predict a quality score that agrees
with human judgement. This is the no-reference image quality assessment (NR-IQA) problem.
Full-reference metrics such as PSNR and SSIM were designed for controlled pipelines where
the original exists; in real usage (camera capture, streaming, social uploads, AI-generated
or AI-enhanced content) it usually does not, and recent work shows the full-reference
paradigm has structural limits \citep{siniukov2022limitations, chen2025afine}. The project
directly extends the course's discussion of objective quality metrics beyond PSNR/SSIM:
we keep the metrics as evaluation tools and build the learned, reference-free alternative.
Concretely, we design a lightweight model that reads a distorted image alone and predicts
its quality, validated against human mean opinion scores.

\section{Related Work}
\label{sec:related}
NR-IQA research currently follows two recipes: (i) reuse a large pretrained backbone with
a lightweight adaptation -- antonym textual prompts for CLIP \citep{wang2023clipiqa,
liao2026maclip}, or a small regressor on frozen features \citep{agnolucci2024arniqa,
yadav2026vlfoundations} -- and (ii) design explicitly lightweight end-to-end models for
deployment \citep{jamshidi2024lariqa}. Surveys of the field \citep{ma2025survey} and the
newest full-reference work \citep{chen2025afine} agree that generalization, not in-dataset
accuracy, is the open problem. Our project takes the frozen-feature recipe and pushes it
to the smallest compute budget: unlike prompt-based methods we use no language component,
and unlike \citet{jamshidi2024lariqa} we do not train a full network -- we freeze a small
backbone and train only a lightweight head on KADID-10k, then check generalization on
TID2013. Table~\ref{tab:refs} summarizes the references that shape this plan.

\begin{table}[ht]
\centering
\caption{Key references.}
\label{tab:refs}
\begin{tabular}{@{}p{0.30\linewidth}p{0.10\linewidth}p{0.52\linewidth}@{}}
\toprule
\textbf{Reference} & \textbf{Year} & \textbf{Relevance to this project} \\
\midrule
\citep{zhang2018lpips} & 2018 & Learned full-reference perceptual metric; our strongest baseline to approach without the reference \\
\citep{ding2022dists}  & 2022 & Structure+texture full-reference metric; optional baseline; texture sensitivity \\
\citep{siniukov2022limitations} & 2022 & Shows differentiable FR metrics can be gamed; motivates the NR framing \\
\citep{wang2023clipiqa} & 2023 & Frozen vision-language prompts for NR-IQA; source of the frozen-backbone philosophy \\
\citep{agnolucci2024arniqa} & 2024 & Frozen encoder + linear quality regressor; closest published design to ours \\
\citep{jamshidi2024lariqa} & 2024 & Lightweight mobile model; uses the same KADID$\rightarrow$TID2013 protocol as our stretch goal \\
\citep{chen2025afine} & 2025 & Evidence that the full-reference paradigm (our baselines' paradigm) has limits; motivation \\
\citep{ma2025survey} & 2025 & Recent survey; used to structure this review and position branches \\
\citep{lao2025spips} & 2025 & Hybrid semantic + low-level perceptual metric; recent reference point \\
\citep{yadav2026vlfoundations} & 2026 & Controlled frozen-backbone comparison; small-head design factors matter \\
\citep{liao2026maclip} & 2026 & Training-free CLIP magnitude-based metric; freshest method reference point \\
\bottomrule
\end{tabular}
\end{table}

\section{Dataset \& Baseline Plan}
\label{sec:data}
\textbf{Dataset.} KADID-10k \citep{lin2019kadid}: 81 pristine images, each degraded by 25
distortion types in 5 levels (10,125 distorted images), with differential mean opinion
scores (DMOS, range 1--5, higher is better) from a crowdsourced study. The database is
freely available for research use and is cited as required. We use the full 10,125 images
and split \emph{by reference image} (65 train / 16 test references, fixed seed 42) so no
content appears in both splits; this prevents the leakage that an image-level split would
cause. TID2013 (3,000 images, 24 distortions) is used only at M3 for a cross-dataset
generalization check. For reproducibility we record the exact archive we downloaded
(byte size 3,067,408,471; SHA-256 of the label file verified against the authors'
mirror) and the exact split seed.

\textbf{Baselines.} PSNR and SSIM (classical, computed on RGB to remain sensitive to the
colour-distortion family in KADID) and the pretrained LPIPS metric \citep{zhang2018lpips}
with the official AlexNet backbone. Correlation is reported per IQA convention: SROCC
(rank correlation, no fitting) and PLCC/RMSE after a 4-parameter logistic mapping, because
metric outputs and DMOS are on different scales. Table~\ref{tab:prelim} reports the
preliminary result on the 2,000-image test split; these numbers are the target our
no-reference model must approach using the distorted image alone.

\begin{table}[ht]
\centering
\caption{Dataset summary.}
\label{tab:data}
\begin{tabular}{@{}p{0.15\linewidth}p{0.47\linewidth}p{0.31\linewidth}@{}}
\toprule
\textbf{Dataset} & \textbf{Subset used} & \textbf{Link} \\
\midrule
KADID-10k & all 10,125 distorted + 81 refs; 65/16 content split (seed 42) & \url{http://database.mmsp-kn.de/kadid-10k-database.html} \\
TID2013 & 3,000 distorted + 25 refs; cross-dataset test only (M3) & \url{http://www.ponomarenko.info/tid2013.htm} \\
\bottomrule
\end{tabular}
\end{table}

\begin{table}[ht]
\centering
\caption{Preliminary baseline correlations on the M1 test split (2,000 distorted
images). All metrics are full-reference and were computed by our own pipeline (code in
the repository).}
\label{tab:prelim}
\begin{tabular}{@{}lccc@{}}
\toprule
\textbf{Metric} & \textbf{SROCC} & \textbf{PLCC} & \textbf{RMSE} \\
\midrule
PSNR  & 0.667 & 0.670 & 0.803 \\
SSIM  & 0.627 & 0.616 & 0.852 \\
LPIPS \citep{zhang2018lpips} & 0.808 & 0.807 & 0.639 \\
\bottomrule
\end{tabular}
\end{table}

\section{Proposed Method}
\label{sec:method}
Our pipeline (to be implemented at M2):
\begin{enumerate}[leftmargin=*]
  \item \textbf{Input:} the distorted image alone, resized to a fixed working
        resolution, with horizontal-flip augmentation.
  \item \textbf{Frozen backbone:} a small ImageNet-pretrained network
        (\todo{primary candidate: MobileNetV3-Large; alternatives for the ablation:
        EfficientNet-B0, a small self-supervised ViT}) that converts the image into a
        pooled feature vector. The backbone is \emph{never} fine-tuned; features can be
        precomputed once per image.
  \item \textbf{Lightweight head:} a 2-layer MLP (with the activation choice informed by
        \citet{yadav2026vlfoundations}) that maps features to one scalar quality score.
        Only these parameters are trained, keeping the model well within the course's
        compute budget and preventing overfitting on 8,125 training images.
  \item \textbf{Training:} regression against DMOS (mean squared error; a pairwise
        ranking term is part of the stretch goal), Adam-family optimizer, early stopping,
        fixed seed 42 for every headline result.
  \item \textbf{Evaluation:} the same content split and the same SROCC/PLCC/RMSE protocol
        as the baselines in Table~\ref{tab:prelim}, plus (stretch) a KADID$\rightarrow$TID2013
        cross-dataset check.
\end{enumerate}
The novelty relative to the cited work is not a new architecture but a rigorously validated,
budget-conscious instantiation: the smallest frozen backbone whose features still support
quality prediction, trained with a protocol that controls leakage, matches scales via the
logistic mapping, and takes generalization seriously.

\section{Feasibility \& Compute Plan}
\label{sec:feasibility}
Compute tier: \textbf{CPU-OK / GPU-light} (per the catalog entry). The full pipeline runs
on the course reference machine (8-core CPU, 16\,GB RAM) and is additionally accelerated
on a free Colab T4 GPU as a supplement. Expected costs: feature precomputation for all
10,125 images takes minutes on GPU / a few hours on CPU and is done once; head training
takes seconds to minutes; the baseline table in Table~\ref{tab:prelim} already ran
end-to-end on the free Colab tier in under ten minutes. All artifacts (checkpoints,
results) are stored both in the repository and in a private drive, and every experiment is
idempotent and seed-fixed, so no result depends on a Colab session surviving.

\subsection{Minimum Viable Deliverable (MVP)}
A frozen small backbone plus a trained lightweight head that predicts DMOS on the M1 test
split, reporting SROCC/PLCC/RMSE in the same table as the baselines -- the smallest version
of the project that still produces a valid, reportable comparison against the baseline.

\subsection{Stretch Goal}
Train on KADID-10k and evaluate on TID2013 (cross-dataset generalization), plus one
ablation isolating the main design choice (backbone size or head capacity, optionally the
ranking-loss term), and an optional comparison against a training-free 2026 reference
point \citep{liao2026maclip} using its released code.

\section{Risk Assessment}
\label{sec:risk}
\begin{itemize}[leftmargin=*]
  \item \textbf{Overfitting on 8,125 training images} (frozen features limit adaptation).
        \emph{Mitigation:} train only a small head, early stopping, flip augmentation, and
        report the validation curve; fallback is a smaller head or feature-space
        regularization.
  \item \textbf{Frozen features may underperform on some distortion families} (e.g.,
        colour distortions the backbone was never optimized for). \emph{Mitigation:}
        per-distortion-type breakdown of correlations, and as a fallback a
        trained-from-scratch compact CNN with stronger augmentation.
  \item \textbf{Colab session limits and environment drift.} \emph{Mitigation:} pinned
        environment file, seed-fixed idempotent scripts, checkpoints and results stored in
        the repository and a private drive, CPU fallback demonstrated.
  \item \textbf{Dataset licensing/ethics.} \emph{Mitigation:} KADID-10k is research-use;
        we cite it, never redistribute the images in the repository, and store only code
        and derived numeric results.
\end{itemize}

\section{Work Breakdown}
\label{sec:workbreakdown}
\begin{table}[ht]
\centering
\caption{Milestone task allocation (see also the project catalog entry for this project).}
\label{tab:tasks}
\begin{tabular}{@{}p{0.12\linewidth}p{0.26\linewidth}p{0.26\linewidth}p{0.26\linewidth}@{}}
\toprule
\textbf{Milestone} & \textbf{Student A} & \textbf{Student B} & \textbf{Student C} \\
\midrule
M1 (Wk 3) & Dataset pipeline + integrity checks; baseline correlation table (done); literature table & Scope frozen-backbone + head design; feasibility test & Evaluation protocol (content split, SROCC/PLCC/ RMSE harness, done); repository/CI hygiene \\
M2 (Wk 6) & Full baseline table; per-distortion-type analysis & Trained head, first checkpoint; backbone ablation started & Preliminary comparison table vs.\ baselines; figures \\
M3 (Wk 9) & Final baselines; TID2013 pipeline & Final model + ablations; ranking-loss variant & Cross-dataset table; demo; report figures \\
\bottomrule
\end{tabular}
\end{table}

\bibliographystyle{unsrtnat}
\bibliography{references}

\end{document}
